\scanpage{211}
\hfill -- \mspage{50} -- \hfill IX

apparaît comme un quotient de $\mathbb Q[T]/(T^{p^N}-1)$, qui est un composé direct de corps extensions finies de $\mathbb Q$. Il en est donc de même de $E_{\mathbb Q}$. Considérant les images de $E,u,v$ dans les corps composants $K_i$ de $E_{\mathbb Q}$, on est donc ramené à voir que si $K=\mathbb Q(u)$ est un corps engendré par une racine primitive $p^n$-ème de l'unité $u$, et si $n$ est un entier $\geq2$ tel que $u-1\equiv0\ \operatorname{mod}n$ (dans l'anneau des entiers de $K$), alors on a $u=1$, ou $u=-1$ et $n=2$. Or ceci résulte aussitôt du fait bien connu que si $u\neq1$ i.e. $m\geq1$, l'idéal $(u-1)$ est premier et l'idéal $(p)$ en est une puissance, dont l'exposant est $>1$ sauf si $p=2,m=1$ i.e. si $u=-1$. Du fait que $(u-1)$ est premier et divise $p$ résulte en effet que $n\geq2$ ne peut diviser $u-1$ que si $n=p$, et $p$ ne divise $u-1$ que si $(u-1)=(p)$ i.e. $u=-1$, donc $n=2$.

4.7.4. Démontrons 4.7.2. Il suffit de prouver que $\Gamma$ est commutatif. Or si $x,y\in\Gamma$, on aura $(xy)(xy)=1$, ce qui peut s'écrire aussi, puisque $x^2=1,y^2=1$, $xy=(xy)^{-1}=y^{-1}x^{-1}=yx$, cqfd.

\uline{Remarque} 4.7.5. La démonstration donnée 4.7.3, jointe à 4.7.2, nous montrent que si $G$ est un Groupe commutatif divisible d'un Topos, $n\geq2$ un entier, et $\Gamma$ un groupe fini, opérant fidèlement sur $G$ en induisant l'identité sur ${}_nG$, alors $\Gamma$ est réduit à l'identité si $n\geq3$, et est isomorphe à un groupe $(\mathbb Z/2\mathbb Z)^r$ si $n=2$.
\margin{368}
\margin{368}
\margin{210}
